
\subsection{典范模型的构造}
\label{subsec:5-4-4}

现在证明前面暂缓的定理 5.4.2 和 5.4.18。
下面的证明与文献 [KoSo01] 的证明相似，
也与本书 2000 年 12 月预印本 [FOOO00]
中仓西映射的构造相似。

\begin{proof}[定理 5.4.2 的证明]
首先定义格林算子与
霍奇–Kodaira 分解定理的组合对应物。
对每个 $k$，固定子空间
\[
\mathcal H^k\subset
\ker\overline{\mathfrak m}_1\cap\overline C^k,
\]
使自然投影
$\ker\overline{\mathfrak m}_1\cap\overline C^k
\to H^k(\overline C;\overline{\mathfrak m}_1)$
限制到 $\mathcal H^k$ 后成为同构。
由于假设
$H^k(\overline C;\overline{\mathfrak m}_1)$
是自由 $R$-模，且 $R$ 是域，
这样的 $\mathcal H^k$ 存在。
取投影
$\ker\overline{\mathfrak m}_1\cap\overline C^k
\to\mathcal H^k$，
使其核恰为
$\operatorname{Im}\overline{\mathfrak m}_1
\cap\overline C^k$，
再任意延拓为
$\overline C^k\to\mathcal H^k$。
将相应的
$\overline C^k\to\overline C^k$ 投影记为
$\Pi_{H,k}$，其像为 $\mathcal H^k$；
并令 $\Pi_H=\sum_k\Pi_{H,k}$。

\begin{lemma}{5.4.28}\label{lem:5-4-28}
存在一列 $R$-模同态
$\mathfrak G_k\colon
\overline C^k\to\overline C^{k-1}$，
满足
\[
1-\Pi_{H,k}
=-\overline{\mathfrak m}_1\circ\mathfrak G_k
-\mathfrak G_{k+1}\circ\overline{\mathfrak m}_1,
\qquad
\mathfrak G_k\circ\mathfrak G_{k+1}=0.
\]
\end{lemma}
\begin{proof}
可取 $F_k\subset\overline C^k$，使
$F_k\oplus\operatorname{Im}\overline{\mathfrak m}_1
=\ker\Pi_{H,k}$。
容易看出，$\overline{\mathfrak m}_1$
诱导同构
$F_k\cong\operatorname{Im}\overline{\mathfrak m}_1
\cap\overline C^{k+1}$。
令 $-\mathfrak G_{k+1}$ 在
$\operatorname{Im}\overline{\mathfrak m}_1$
上为此同构的逆，
在 $F_{k+1}+\mathcal H^{k+1}$ 上为零。
由此立即得到断言。
\end{proof}

\begin{remark}{5.4.29}\label{rem:5-4-29}
仅在本引理和下文引理 5.4.51 中
使用 $R$ 为域的假设。
\end{remark}

令 $G\subset\mathbb R_{\geq0}\times2\mathbb Z$
如条件 3.1.6，
并把 $G$ 到 $\mathbb R_{\geq0}$
的投影按严格递增次序写成
$\{\lambda_i\mid i=0,1,2,\ldots\}$，
其中 $\lambda_0=0$。
假设所有滤过 $A_\infty$ 代数及模都 $G$-有隙。
于是
\begin{equation}
\mathfrak m_k
=\sum_i\mathfrak m_{k,i}T^{\lambda_i},
\tag{5.4.30.1}\label{eq:5-4-30-1}
\end{equation}
其中 $\mathfrak m_{k,i}$ 不含 $T$，即
\[
\mathfrak m_{k,i}\colon
B_k\overline C[1]\otimes R[e,e^{-1}]
\longrightarrow\overline C[1]\otimes R[e,e^{-1}].
\]
记
\begin{equation}
\mathfrak m_{k,i}^{\circ}
=\mathfrak m_{k,i}T^{\lambda_i}.
\tag{5.4.30.2}\label{eq:5-4-30-2}
\end{equation}
设 $C_{\mathrm{can}}=\mathcal H\otimes\Lambda_{0,\mathrm{nov}}$。
下面将在其上定义滤过 $A_\infty$ 结构
$\mathfrak m'\,(=\mathfrak m_{\mathrm{can}})$；
其运算 $\mathfrak m'_k$ 是在树上求和得到的。

先把 $\Pi_{H,k}$ 和 $\mathfrak G_k$
自然延拓为
\[
\begin{aligned}
\Pi_{H,k}&\colon
\overline C^k\otimes R[e,e^{-1}]
\longrightarrow\overline C^k\otimes R[e,e^{-1}],\\
\mathfrak G_k&\colon
\overline C^k\otimes R[e,e^{-1}]
\longrightarrow\overline C^{k-1}\otimes R[e,e^{-1}],
\end{aligned}
\]
并沿用原记号。

现定义树集 $\mathcal G_k^+$。
它是第 3.4 节引入的带状树模空间
$\mathcal G_k$ 的一个小变形：
$\mathcal G_k^+$ 的元素为五元组
$\Gamma=(T,i,v_0,V_{\mathrm{tad}},\eta)$，
满足以下条件：
\begin{enumerate}
\item[\textup{(5.4.31.1)}]\label{cond:5-4-31-1}
$(T,i,v_0)$ 满足条件 (3.4.1.3)，
但不一定满足 (3.4.1.1)、(3.4.1.2)；
即 $T$ 可以有二价顶点，
也可以有一价内顶点。
\item[\textup{(5.4.31.2)}]\label{cond:5-4-31-2}
$V_{\mathrm{tad}}$ 是一价顶点集的一个子集，
可以为空；全体一价顶点分解为
$C^0_{\mathrm{ext}}(T)$ 与
$V_{\mathrm{tad}}$ 的不交并，
且 $v_0\in C^0_{\mathrm{ext}}(T)$。
置 $C^0_{\mathrm{int}}(T)=
C^0(T)\setminus C^0_{\mathrm{ext}}(T)$，
并令 $k=\#C^0_{\mathrm{ext}}(T)$。
又有
$C^0_{\mathrm{ext}}(T)
=T\cap i^{-1}(\partial D^2)$。
\item[\textup{(5.4.31.3)}]\label{cond:5-4-31-3}
$\eta\colon C^0_{\mathrm{int}}(T)
\to\{0,1,2,\ldots\}$；
若顶点 $v$ 仅有一条或两条边，
则要求 $\eta(v)>0$。
\end{enumerate}
$V_{\mathrm{tad}}$ 中的元素称为
\liftIndex{I214}
蝌蚪端点（tad pole），并视作内顶点。
换言之，选择 $V_{\mathrm{tad}}$
即可指定 $C^0_{\mathrm{ext}}(T)$
和 $C^0_{\mathrm{int}}(T)$。

对于
$\Gamma=(T,i,v_0,V_{\mathrm{tad}},\eta)
\in\mathcal G_{k+1}^+$，
将关联两个同态：
\begin{equation}
\mathfrak m_\Gamma\colon
B_k\mathcal H[1]\otimes R[e,e^{-1}]
\longrightarrow\mathcal H[1]\otimes R[e,e^{-1}],
\quad\deg\mathfrak m_\Gamma=1;
\tag{5.4.32.1}\label{eq:5-4-32-1}
\end{equation}
\begin{equation}
\mathfrak f_\Gamma\colon
B_k\mathcal H[1]\otimes R[e,e^{-1}]
\longrightarrow\overline C[1]\otimes R[e,e^{-1}],
\quad\deg\mathfrak f_\Gamma=0.
\tag{5.4.32.2}\label{eq:5-4-32-2}
\end{equation}
以下按树的次序归纳定义这些算子。
对 $\Gamma\in\mathcal G_{k+1}^+$，
定义其能量
\[
E(\Gamma)=
\sum_{v\in C^0_{\mathrm{int}}(T)}\lambda_{\eta(v)}.
\]
若 $\Gamma'\in\mathcal G_{k'+1}^+$，
称 $\Gamma'>\Gamma$，当且仅当以下
两种情况之一成立：
\begin{enumerate}
\item[\textup{(5.4.33.1)}]\label{cond:5-4-33-1}
$E(\Gamma')>E(\Gamma)$；
\item[\textup{(5.4.33.2)}]\label{cond:5-4-33-2}
$E(\Gamma')=E(\Gamma)$ 且 $k'>k$。
\end{enumerate}

先考虑没有内顶点的情形。
此时仅有一棵树 $\Gamma_0$，令
\[
\mathfrak m_{\Gamma_0}=0,
\qquad
\mathfrak f_{\Gamma_0}=
\mathrm{incl}\colon\mathcal H[1]\to\overline C[1].
\]

现在考虑仅有一个内顶点的树。
对每个 $k\geq0$ 及 $\eta\geq0$，
只要 $(k,\eta)\ne(0,0),(1,0)$，
就有唯一一棵这样的
$\Gamma_{k+1,\eta}\in\mathcal G_{k+1}^+$。
这里 $\eta$ 视为内顶点 $v$ 处的数
$\eta(v)$；$k=0$ 时，$v$ 是蝌蚪端点。
图 5.4.1 中，$v_0$ 放在顶部，
输入顶点按逆时针次序排列。
\translatornote{原书此处把参数写成
$k=1,2,\ldots$，却又讨论 $k=0$
并画出 $\Gamma_{1,\eta}$；
且把一、二价内顶点的正能量条件
误说成 $k=1,2$。
依条件 (5.4.31.3)，
应纳入 $k=0$，正能量限制对应
$k=0,1$；见疑误 A-0326。}
\translatornote{原书式 (5.4.34)
第一行把定义域写为
$B_k\overline C[1]$，
与 (5.4.32.1) 的树运算输入不符；
这里改为 $B_k\mathcal H[1]$，
见疑误 A-0325。}
定义
\begin{equation}
\left\{\begin{aligned}
\mathfrak m_{\Gamma_{k+1,\eta}}
&=\Pi_H\circ\mathfrak m_{k,\eta(v)}
\colon B_k\mathcal H[1]\otimes R[e,e^{-1}]
\longrightarrow\mathcal H[1]\otimes R[e,e^{-1}],\\
\mathfrak f_{\Gamma_{k+1,\eta}}
&=\mathfrak G\circ\mathfrak m_{k,\eta(v)}
\colon B_k\mathcal H[1]\otimes R[e,e^{-1}]
\longrightarrow\overline C[1]\otimes R[e,e^{-1}].
\end{aligned}\right.
\tag{5.4.34}\label{eq:5-4-34}
\end{equation}
右端的 $\mathfrak m_{k,\eta(v)}$
来自 (5.4.30.1)。
当 $k=0,1$ 时，
由条件 (5.4.31.3) 有 $\eta(v)>0$；
因而典范模型满足
$\overline{\mathfrak m}'_0
=\overline{\mathfrak m}'_1=0$。

\begin{figure}[!ht]
\centering
\setlength{\unitlength}{0.8mm}
\begin{picture}(155,52)(0,-8)
 \put(13,36){\line(0,-1){23}}
 \put(13,36){\circle*{1.5}}
 \put(13,13){\circle*{1.5}}
 \put(43,36){\line(0,-1){27}}
 \put(43,36){\circle*{1.5}}
 \put(43,22){\circle*{1.5}}
 \put(43,9){\circle*{1.5}}
 \put(78,36){\line(0,-1){17}}
 \put(78,19){\line(-1,-1){12}}
 \put(78,19){\line(1,-1){12}}
 \put(78,36){\circle*{1.5}}
 \put(78,19){\circle*{1.5}}
 \put(66,7){\circle*{1.5}}
 \put(90,7){\circle*{1.5}}
 \put(126,36){\line(0,-1){17}}
 \put(126,19){\line(-1,-1){17}}
 \put(126,19){\line(1,-1){17}}
 \put(126,36){\circle*{1.5}}
 \put(126,19){\circle*{1.5}}
 \put(109,2){\circle*{1.5}}
 \put(143,2){\circle*{1.5}}
 \put(118,6){$\cdots$}
 \put(8,39){$\scriptstyle v_0$}
 \put(39,39){$\scriptstyle v_0$}
 \put(74,39){$\scriptstyle v_0$}
 \put(122,39){$\scriptstyle v_0$}
 \put(14,19){$\scriptstyle v$}
 \put(44,23){$\scriptstyle v$}
 \put(80,19){$\scriptstyle v$}
 \put(128,19){$\scriptstyle v$}
 \put(2,1){$\scriptstyle\eta(v)>0$}
 \put(35,1){$\scriptstyle\eta(v)>0$}
 \put(5,-7){$\Gamma_{1,\eta}$}
 \put(36,-7){$\Gamma_{2,\eta}$}
 \put(70,-7){$\Gamma_{3,\eta}$}
 \put(116,-7){$\Gamma_{k+1,\eta}$}
\end{picture}
\par\smallskip
\renewcommand{\thefigure}{5.4.1}
\refstepcounter{figure}
{\normalfont\bfseries 图~\thefigure}\label{fig:5-4-1}
\end{figure}

若 $\Gamma=(T,i,v_0,V_{\mathrm{tad}},\eta)
\in\mathcal G_{k+1}^+$ 有多个内顶点，
令 $v_1$ 为与 $v_0$ 以一条边相连的
唯一顶点。
设 $v_1$ 有 $\ell+1$ 条边，
按嵌入 $i\colon T\to D^2$
给出的逆时针顺序将它们记作
$e_0,\ldots,e_\ell$，
从连接 $v_1$ 与 $v_0$ 的 $e_0$ 开始。
在 $e_1,\ldots,e_\ell$ 的中点切开 $T$，
得到
$\Gamma_j\in\mathcal G^+_{k_j+1}$
（$1\leq j\leq\ell$），
其中 $k=\sum_j k_j$；见图 5.4.2。

\begin{figure}[!ht]
\centering
\setlength{\unitlength}{0.85mm}
\begin{picture}(112,55)(0,-5)
 \put(56,46){\circle*{1.5}}
 \put(56,46){\line(0,-1){19}}
 \put(56,27){\circle*{1.5}}
 \put(56,27){\line(-2,-1){30}}
 \put(56,27){\line(-1,-2){9}}
 \put(56,27){\line(1,-2){9}}
 \put(56,27){\line(2,-1){30}}
 \put(30,17){$\cdots$}
 \put(74,17){$\cdots$}
 \put(58,46){$\scriptstyle v_0$}
 \put(58,27){$\scriptstyle v_1$}
 \put(26,24){$\scriptstyle e_1$}
 \put(43,18){$\scriptstyle e_2$}
 \put(66,18){$\scriptstyle e_{\ell-1}$}
 \put(86,24){$\scriptstyle e_\ell$}
 \put(15,6){\oval(24,17)}
 \put(43,1){\oval(17,18)}
 \put(69,1){\oval(17,18)}
 \put(97,6){\oval(24,17)}
 \put(7,5){$\Gamma_1$}
 \put(36,0){$\Gamma_2$}
 \put(61,0){$\Gamma_{\ell-1}$}
 \put(90,5){$\Gamma_\ell$}
\end{picture}
\par\smallskip
\renewcommand{\thefigure}{5.4.2}
\refstepcounter{figure}
{\normalfont\bfseries 图~\thefigure}\label{fig:5-4-2}
\end{figure}

\begin{lemma}{5.4.35}\label{lem:5-4-35}
对每个 $j$，有 $\Gamma_j<\Gamma$。
\end{lemma}
\begin{proof}
由于
$\sum_jE(\Gamma_j)+\lambda_{\eta(v_1)}
=E(\Gamma)$，
当 $\lambda_{\eta(v_1)}>0$ 时结论显然。
设 $\lambda_{\eta(v_1)}=0$，
于是 $\ell\geq2$。
若 $E(\Gamma_j)<E(\Gamma)$，
仍有 $\Gamma_j<\Gamma$。
若 $E(\Gamma_j)=E(\Gamma)$，
则 $\Gamma\setminus\Gamma_j$
上没有蝌蚪端点；
由 $\ell\geq2$ 可知其上有外顶点，
所以 $k_j<k$。
\end{proof}

对 $x\in B_k\mathcal H[1]\otimes R[e,e^{-1}]$，
写
\[
\Delta^{\ell-1}x
=\sum_a x_a^{(1)}\otimes\cdots\otimes x_a^{(\ell)}.
\]
归纳地定义
\begin{equation}
\left\{\begin{aligned}
\mathfrak m_\Gamma(x)
&=\sum_a\Pi_H\circ
\mathfrak m_{\ell,\eta(v_1)}
\bigl(\mathfrak f_{\Gamma_1}(x_a^{(1)})
\otimes\cdots\otimes
\mathfrak f_{\Gamma_\ell}(x_a^{(\ell)})\bigr),\\
\mathfrak f_\Gamma(x)
&=\sum_a\mathfrak G\circ
\mathfrak m_{\ell,\eta(v_1)}
\bigl(\mathfrak f_{\Gamma_1}(x_a^{(1)})
\otimes\cdots\otimes
\mathfrak f_{\Gamma_\ell}(x_a^{(\ell)})\bigr).
\end{aligned}\right.
\tag{5.4.36}\label{eq:5-4-36}
\end{equation}
若 $x_a^{(j)}$ 不属于
$B_{k_j}\mathcal H[1]\otimes R[e,e^{-1}]$，
则约定 $\mathfrak f_{\Gamma_j}(x_a^{(j)})=0$。
引理 5.4.35 保证右端已定义。

\begin{definition}{5.4.37}\label{def:5-4-37}
令
\[
\mathfrak m'_k
=\sum_{\Gamma\in\mathcal G^+_{k+1}}
T^{E(\Gamma)}\mathfrak m_\Gamma,
\qquad
\mathfrak f_k
=\sum_{\Gamma\in\mathcal G^+_{k+1}}
T^{E(\Gamma)}\mathfrak f_\Gamma.
\]
\end{definition}

下面证明 $\mathfrak m'_k$ 在
$C_{\mathrm{can}}$ 上定义
$\overline{\mathfrak m}'_1=0$
的滤过 $A_\infty$ 结构，
且 $\mathfrak f_k$ 定义从
$(C_{\mathrm{can}},\mathfrak m')$
到 $(C,\mathfrak m)$ 的
$A_\infty$ 同态。
令 $\widehat{\mathfrak f}$
为 $\mathfrak f_k$ 诱导的余代数同态，
$\widehat{\mathfrak m}'_k$
为 $\mathfrak m'_k$ 诱导的余导子，
并置
$\widehat d'=\sum_k\widehat{\mathfrak m}'_k$。

\begin{lemma}{5.4.38}\label{lem:5-4-38}
$\widehat{\mathfrak f}\circ\widehat d'
=\widehat d\circ\widehat{\mathfrak f}$。
\end{lemma}
\begin{proof}
在 $\{(k,i)\mid k,i\in\mathbb Z_{\geq0}\}$
上定义次序：
$(k',i')<(k,i)$ 当且仅当
$i'<i$，或 $i'=i$ 且 $k'<k$。
将按此顺序归纳证明
\begin{equation}
\widehat{\mathfrak f}\circ\widehat d'
\equiv\widehat d\circ\widehat{\mathfrak f}
\pmod{T^{\lambda_i}}
\quad\text{在 }B_kC_{\mathrm{can}}[1]\text{ 上}
\tag{5.4.39.$(k,i)$}\label{eq:5-4-39}
\end{equation}
起始情形 $i=0$ 由
$\mathfrak m_{\Gamma_0}=0$
及 $\mathfrak f_{\Gamma_0}$ 为包含映射即得。
假设 (5.4.39.$(k',i')$)
对所有 $(k',i')<(k,i)$ 成立。
由树递归的定义有
\translatornote{原书式 (5.4.40)
第一行漏去 $k=1$ 时
无内顶点树的包含映射贡献。
这里先将该项从 $\mathfrak f_k$
扣除再写递归和式；见疑误 A-0329。}
\begin{equation}
\left\{\begin{aligned}
\mathfrak f_k-\delta_{k,1}\mathfrak f_{\Gamma_0}
&=\sum_{(\ell,j)\ne(1,0)}
\mathfrak G\circ\mathfrak m^\circ_{\ell,j}
\circ\widehat{\mathfrak f}
&&\text{在 }B_kC_{\mathrm{can}}[1]\text{ 上},\\
\mathfrak m'_k
&=\sum_{(\ell,j)\ne(1,0)}
\Pi_H\circ\mathfrak m^\circ_{\ell,j}
\circ\widehat{\mathfrak f}
&&\text{在 }B_kC_{\mathrm{can}}[1]\text{ 上}.
\end{aligned}\right.
\tag{5.4.40}\label{eq:5-4-40}
\end{equation}
这里 $\mathfrak G$ 与 $\Pi_H$
都平凡延拓到 $\Lambda_{0,\mathrm{nov}}$ 上。
在 $B_kC_{\mathrm{can}}[1]$ 上计算得
\begin{equation}
\begin{aligned}
\widehat d\circ\widehat{\mathfrak f}
&=\mathfrak m^\circ_{1,0}\circ\widehat{\mathfrak f}
 +\sum_{(\ell,j)\ne(1,0)}
 \mathfrak m^\circ_{\ell,j}\circ\widehat{\mathfrak f}\\
&=\sum_{\Gamma\in\mathcal G_{k+1}^+}
T^{E(\Gamma)}\mathfrak m^\circ_{1,0}
\circ\mathfrak f_\Gamma
 +\sum_{(\ell,j)\ne(1,0)}
\mathfrak m^\circ_{\ell,j}\circ\widehat{\mathfrak f}\\
&=\sum_{(\ell,j)\ne(1,0)}
\mathfrak m^\circ_{1,0}\circ\mathfrak G
\circ\mathfrak m^\circ_{\ell,j}
\circ\widehat{\mathfrak f}
 +\sum_{(\ell,j)\ne(1,0)}
\mathfrak m^\circ_{\ell,j}\circ\widehat{\mathfrak f}.
\end{aligned}
\tag{5.4.41}\label{eq:5-4-41}
\end{equation}
对第一项应用引理 5.4.28，
其中一项与第二项相消，留下
\begin{equation}
\begin{aligned}
\widehat d\circ\widehat{\mathfrak f}
={}&-\sum_{(\ell,j)\ne(1,0)}
\mathfrak G\circ\mathfrak m^\circ_{1,0}
\circ\mathfrak m^\circ_{\ell,j}
\circ\widehat{\mathfrak f}\\
&+\sum_{(\ell,j)\ne(1,0)}
\Pi_H\circ\mathfrak m^\circ_{\ell,j}
\circ\widehat{\mathfrak f}.
\end{aligned}
\tag{5.4.42}\label{eq:5-4-42}
\end{equation}
第二项就是
$\mathfrak m'_k=\widehat d'
=\mathfrak f_{\Gamma_0}\circ\widehat d'$
在 $B_kC_{\mathrm{can}}[1]$ 上的作用。
由 $A_\infty$ 关系，
第一项等于
$\sum_{(\ell,j)\ne(1,0)}
\mathfrak G\circ\mathfrak m^\circ_{\ell,j}
\circ\widehat d\circ\widehat{\mathfrak f}$。
因此
\begin{equation}
\widehat d\circ\widehat{\mathfrak f}
=\sum_{(\ell,j)\ne(1,0)}
\mathfrak G\circ\mathfrak m^\circ_{\ell,j}
\circ\widehat d\circ\widehat{\mathfrak f}
+\mathfrak f_{\Gamma_0}\circ\widehat d'.
\tag{5.4.43}\label{eq:5-4-43}
\end{equation}
只须证明 (5.4.43) 的第一项模
$T^{\lambda_i}$ 同余于
$\sum_{(\ell,j)\ne(1,0)}
\mathfrak G\circ\mathfrak m^\circ_{\ell,j}
\circ\widehat{\mathfrak f}\circ\widehat d'$。

先考虑 $\ell=0,1$。
由于 $\mathfrak m_{0,0}=0$，
此时 $j>0$，故
$\mathfrak m^\circ_{\ell,j}
\equiv0\pmod{T^{\lambda_1}}$。
归纳假设给出
$\widehat{\mathfrak f}\circ\widehat d'
\equiv\widehat d\circ\widehat{\mathfrak f}
\pmod{T^{\lambda_{i-1}}}$，
于是
\begin{equation}
\mathfrak G\circ\mathfrak m^\circ_{\ell,j}
\circ\widehat d\circ\widehat{\mathfrak f}
\equiv
\mathfrak G\circ\mathfrak m^\circ_{\ell,j}
\circ\widehat{\mathfrak f}\circ\widehat d'
\pmod{T^{\lambda_i}}
\tag{5.4.44}\label{eq:5-4-44}
\end{equation}
对 $\ell=0,1$ 且 $(\ell,j)\ne(1,0)$ 成立。
再考虑 $\ell\geq2$。
令 $x\in B_kC_{\mathrm{can}}[1]$，
将 $(\widehat{\mathfrak f}\circ\widehat d')(x)$
的 $B_\ell C[1]$ 分量记为
$(\widehat{\mathfrak f}\circ\widehat d')_\ell(x)$。
按约定 (Conv.3) 的 (1.35)，有
\begin{equation}
\begin{aligned}
(\widehat{\mathfrak f}\circ\widehat d')_\ell(x)
={}&\sum_a\sum_j
(-1)^{\deg' x_{1;a}+\cdots+\deg' x_{j-1;a}}\\
&\quad\mathfrak f(x_{1;a})\otimes\cdots\otimes
(\mathfrak f\circ\widehat d')(x_{j;a})
\otimes\cdots\otimes\mathfrak f(x_{\ell;a}).
\end{aligned}
\tag{5.4.45}\label{eq:5-4-45}
\end{equation}
定义 $k_{j;a}$ 使
$x_{j;a}\in B_{k_{j;a}}C_{\mathrm{can}}[1]$，
于是
\begin{equation}
\sum_j k_{j;a}=k.
\tag{5.4.46}\label{eq:5-4-46}
\end{equation}
若每个 $k_{j;a}<k$，
归纳假设给出
\begin{equation}
\begin{aligned}
&\mathfrak f(x_{1;a})\otimes\cdots\otimes
(\mathfrak f\circ\widehat d')(x_{j;a})
\otimes\cdots\otimes\mathfrak f(x_{\ell;a})\\
&\quad\equiv
\mathfrak f(x_{1;a})\otimes\cdots\otimes
(\mathfrak m\circ\widehat{\mathfrak f})(x_{j;a})
\otimes\cdots\otimes\mathfrak f(x_{\ell;a})
\pmod{T^{\lambda_i}}.
\end{aligned}
\tag{5.4.47}\label{eq:5-4-47}
\end{equation}
最后设某个 $k_{j;a}=k$。
由 (5.4.46)，
所有 $j'\ne j$ 均有 $k_{j';a}=0$。
此时 (5.4.47) 左端的张量项为
\[
\mathfrak f(1)^{\otimes(j-1)}
\otimes(\mathfrak f\circ\widehat d')(x)
\otimes\mathfrak f(1)^{\otimes(\ell-j)}.
\]
由假设
$\mathfrak f(1)\equiv0\pmod{T^{\lambda_1}}$。
由于 $\ell\geq2$，
再用归纳假设
$(\mathfrak f\circ\widehat d')(x)
\equiv(\mathfrak m\circ\widehat{\mathfrak f})(x)
\pmod{T^{\lambda_{i-1}}}$，
可知 (5.4.47) 在这一情形也成立。

于是 (5.4.44)、(5.4.47) 与 (5.4.36)
推出
\[
\begin{aligned}
&\sum_{(\ell,j)\ne(1,0)}
\mathfrak G\circ\mathfrak m^\circ_{\ell,j}
\circ\widehat d\circ\widehat{\mathfrak f}\\
&\quad\equiv
\left(\sum_{(\ell,j)\ne(1,0)}
\mathfrak G\circ\mathfrak m^\circ_{\ell,j}
\circ\widehat{\mathfrak f}\right)
\circ\widehat d'\\
&\quad\equiv
(\widehat{\mathfrak f}
-\mathfrak f_{\Gamma_0})\circ\widehat d'
\pmod{T^{\lambda_i}}.
\end{aligned}
\]
结合 (5.4.43) 即完成引理的证明。
\end{proof}

由于 $\mathfrak f_{\Gamma_0}$ 是包含映射，
而 $\widehat{\mathfrak f}$ 模
$\Lambda^+_{0,\mathrm{nov}}$ 与该包含映射相同，
故 $\widehat{\mathfrak f}$ 单射。
于是引理 5.4.38 蕴涵
$\widehat d'\circ\widehat d'=0$，
进而 $\mathfrak f$ 是 $A_\infty$ 同态。
从构造易见，$\mathfrak f$
在上同调上诱导同构。
由定理 4.2.45 即得定理 5.4.2。
\end{proof}

\begin{proof}[定理 5.4.2$'$ 中幺元断言的证明]
最后证明关于幺元的断言。
若 $(C,\mathfrak m)$ 同伦含幺，
则按定义 3.3.2 从
$(C^+,\mathfrak m^+)$ 开始上述构造。
故只需讨论 $(C,\mathfrak m)$ 含幺的情形。
可选 $\mathcal H$ 使它包含单位元 $\mathbf e$。
下面将按 (5.4.33) 给出的树次序归纳证明：
\begin{align}
\mathfrak f_\Gamma
(y_1\otimes\cdots\otimes y_\ell
\otimes\mathbf e\otimes
z_1\otimes\cdots\otimes z_{k-\ell})
&=0\quad(k>0),
\tag{5.4.48.1}\label{eq:5-4-48-1}\\
\mathfrak f_\Gamma(\mathbf e)
&=0\quad(E(\Gamma)>0),
\tag{5.4.48.2}\label{eq:5-4-48-2}\\
\mathfrak f_{\Gamma_0}(\mathbf e)
&=\mathbf e\quad(E(\Gamma_0)=0).
\tag{5.4.48.3}\label{eq:5-4-48-3}
\end{align}
沿用 (5.4.36) 的记号。
考虑 $x=y\otimes\mathbf e\otimes z$。
由归纳假设，
$\mathfrak f_\Gamma(x)$ 仅可能为
\translatornote{原书式 (5.4.49)
把 $\mathfrak m$ 的第二下标写作
$\lambda_{\eta(v)}$；
依 (5.4.30.1)、(5.4.36)
改为整数序号 $\eta(v)$，
见疑误 A-0332。}
\begin{equation}
\begin{aligned}
\mathfrak f_\Gamma(x)
=\sum_{a,b}
(&\mathfrak G\circ\mathfrak m_{\ell,\eta(v)})
\bigl(\mathfrak f_{\Gamma_{1,1}}(y_a^1)
\otimes\cdots\otimes
\mathfrak f_{\Gamma_{1,\alpha}}(y_a^\alpha)\\
&\qquad\otimes\mathbf e\otimes
\mathfrak f_{\Gamma_{2,1}}(z_b^1)
\otimes\cdots\otimes
\mathfrak f_{\Gamma_{2,\beta}}(z_b^\beta)\bigr).
\end{aligned}
\tag{5.4.49}\label{eq:5-4-49}
\end{equation}
因 $\mathbf e$ 是单位元，
除非 $\alpha+\beta=1$ 且
$\lambda_{\eta(v)}=0$，
(5.4.49) 均为零。
在余下的情形，
\begin{equation}
\mathfrak f_\Gamma(x)
=\mathfrak G(\mathfrak f_{\Gamma_1}(y))
\quad\text{或}\quad
\pm\mathfrak G(\mathfrak f_{\Gamma_2}(z)).
\tag{5.4.50}\label{eq:5-4-50}
\end{equation}
先考虑第一种。
若 $k>1$ 或 $E(\Gamma_1)>0$，
则由 (5.4.40)
$\mathfrak f_{\Gamma_1}(y)
\in\operatorname{Im}\mathfrak G$；
由于 $\mathfrak G\circ\mathfrak G=0$，
(5.4.50) 此时为零。
若 $k=1$ 且 $E(\Gamma_1)=0$，
则
$\mathfrak f_{\Gamma_1}(y)=y_1\in\mathcal H$，
故 $\mathfrak G(y_1)=0$。
\translatornote{原书此处误写成
$\mathfrak G(\mathfrak f_{\Gamma_1}(y))=y_1$；
零能量基例给出的是
$\mathfrak f_{\Gamma_1}(y)=y_1$，
见疑误 A-0331。}
第二种情形相同，故 (5.4.48) 得证。
再用 (5.4.48)，
可类似证明 $\mathbf e$ 是典范模型的单位元；
同时它也表明 $\mathfrak f$ 保持单位元。
\end{proof}

下面证明定理 5.4.18。
\begin{proof}[定理 5.4.18 的证明]
先给出引理 5.4.28 的双模对应版本。
令 $D$ 如定理 5.4.18，
并令 $\overline D$ 为无滤过
$A_\infty$ 双模，使
$D=\overline D\otimes\Lambda_{0,\mathrm{nov}}$。
对每个 $k$，取
\[
\mathcal H^k\subset
\ker\overline{\mathfrak n}_{0,0}\cap\overline D^k,
\]
使从
$\ker\overline{\mathfrak n}_{0,0}\cap\overline D^k$
到
$H^k(\overline D;\overline{\mathfrak n}_{0,0})$
的投影限制到 $\mathcal H^k$ 后为同构。
这里用到该上同调为自由模的假设。
\translatornote{原书以下把投影的定义域
和上述上同调对象误写为 $\overline C$；
$\overline{\mathfrak n}_{0,0}$
作用在双模 $\overline D$ 上，
故统一为 $\overline D$，见疑误 A-0327。}
取投影
$\ker\overline{\mathfrak n}_{0,0}\cap\overline D^k
\to\mathcal H^k$，
使其核为
$\operatorname{Im}\overline{\mathfrak n}_{0,0}
\cap\overline D^k$，
再任意延拓到 $\overline D^k\to\mathcal H^k$。
相应的幂等映射
$\overline D^k\to\overline D^k$
记为 $\Pi_{H,k}$，像为 $\mathcal H^k$；
置 $\Pi_H=\sum_k\Pi_{H,k}$。

\translatornote{原书引理 5.4.51
把两个降一次的 $\mathfrak G$
写作 $\mathfrak G_{k+1}\circ\mathfrak G_k$，
型别不合；这里依引理 5.4.28
调整复合次序，见疑误 A-0330。}
\begin{lemma}{5.4.51}\label{lem:5-4-51}
存在一列映射
$\mathfrak G_k\colon
\overline D^k\to\overline D^{k-1}$，
满足
\[
1-\Pi_{H,k}
=-\overline{\mathfrak n}_{0,0}\circ\mathfrak G_k
-\mathfrak G_{k+1}\circ\overline{\mathfrak n}_{0,0},
\qquad
\mathfrak G_k\circ\mathfrak G_{k+1}=0.
\]
\end{lemma}
其证明与引理 5.4.28 相同，从略。
令
\[
\mathfrak n_{k_1,k_0}
=\sum_iT^{\lambda_i}\mathfrak n_{k_1,k_0;i},
\]
其中 $\mathfrak n_{k_1,k_0;i}$ 不含 $T$。
以下借助图 5.4.3 所示费曼图，
构造
$D'=H(\overline D;\overline{\mathfrak n}_{0,0})
\otimes\Lambda_{0,\mathrm{nov}}$
上的 $A_\infty$ 双模结构 $\mathfrak n'$，
以及双模同态 $\varphi\colon D'\to D$。

\begin{figure}[!ht]
\centering
\setlength{\unitlength}{0.85mm}
\begin{picture}(135,93)(0,-7)
 \put(65,84){\circle*{1.6}}
 \put(65,84){\line(0,-1){78}}
 \put(65,6){\circle*{1.6}}
 \put(65,66){\circle*{1.6}}
 \put(65,43){\circle*{1.6}}
 \put(65,20){\circle*{1.6}}
 \put(40,66){\line(1,0){50}}
 \put(40,66){\circle*{1.6}}
 \put(90,66){\circle*{1.6}}
 \put(21,43){\line(1,0){72}}
 \put(21,43){\circle*{1.6}}
 \put(40,43){\circle*{1.6}}
 \put(93,43){\circle*{1.6}}
 \qbezier(93,43)(106,47)(118,50)
 \qbezier(93,43)(106,39)(118,36)
 \put(118,50){\circle*{1.6}}
 \put(118,36){\circle*{1.6}}
 \put(40,20){\line(1,0){72}}
 \put(40,20){\circle*{1.6}}
 \put(90,20){\circle*{1.6}}
 \put(112,20){\circle*{1.6}}
 \put(37,72){$\scriptstyle\nu_1=0$}
 \put(91,72){$\scriptstyle\nu_0=0$}
 \put(14,48){$\scriptstyle\nu_1=1$}
 \put(118,44){$\scriptstyle\nu_0=2$}
 \put(36,14){$\scriptstyle\nu_1=0$}
 \put(113,14){$\scriptstyle\nu_0=1$}
 \put(68,64){$\scriptstyle\mu(3)=2$}
 \put(68,41){$\scriptstyle\mu(2)=0$}
 \put(68,18){$\scriptstyle\mu(1)=2$}
\end{picture}
\par\smallskip
\renewcommand{\thefigure}{5.4.3}
\refstepcounter{figure}
{\normalfont\bfseries 图~\thefigure}\label{fig:5-4-3}
\end{figure}

设 $k$ 为非负整数。
\begin{definition}{5.4.52}\label{def:5-4-52}
定义 $\Theta(k)$ 为全部三元组
$(\nu_1,\mu,\nu_0)$ 构成的集合，
其中
$\nu_i\colon\{1,\ldots,k\}\to\mathbb Z_{\geq0}$
（$i=0,1$），
$\mu\colon\{1,\ldots,k\}\to\mathbb Z_{\geq0}$，
并且
\begin{equation}
\nu_1(j)+\mu(j)+\nu_0(j)>0
\quad\text{对一切 }j=1,\ldots,k.
\tag{5.4.53}\label{eq:5-4-53}
\end{equation}
记
\[
|\nu_1|=\sum_j\nu_1(j),\qquad
|\mu|=\sum_j\lambda_{\mu(j)},\qquad
|\nu_0|=\sum_j\nu_0(j).
\]
\end{definition}

图 5.4.3 对应于
$\nu_1(1)=\nu_1(3)=\nu_0(3)=0$、
$\nu_1(2)=1$、$\nu_0(1)=1$、
$\nu_0(2)=2$，
以及 $\mu(1)=2$、$\mu(2)=0$、
$\mu(3)=2$。
对每个
$(\nu_1,\mu,\nu_0)\in\Theta(k)$，
将归纳定义两个映射
\begin{equation}
\left\{\begin{aligned}
\varphi_{(\nu_1,\mu,\nu_0)}&\colon
B_{|\nu_1|}C_1[1]\otimes D'[1]\otimes
B_{|\nu_0|}C_0[1]\longrightarrow D[1],\\
\mathfrak n'_{(\nu_1,\mu,\nu_0)}&\colon
B_{|\nu_1|}C_1[1]\otimes D'[1]\otimes
B_{|\nu_0|}C_0[1]\longrightarrow D'[1].
\end{aligned}\right.
\tag{5.4.54}\label{eq:5-4-54}
\end{equation}
当 $k=0$ 时，$\Theta(0)$ 仅有
一个元素，记为 $\varnothing$；
令 $\varphi_\varnothing=\mathrm{id}$、
$\mathfrak n'_\varnothing
=\mathfrak n_{0,0;0}|_{\mathcal H}=0$。
\translatornote{原书基例只写
$\mathfrak n'_\varnothing=\mathfrak n_{0,0;0}$；
这里明确将约化微分限制到上同调代表空间
$\mathcal H$，使 (5.4.54) 的型别成立，
见疑误 A-0333。}
现设对 $k-1$ 已定义 (5.4.54)，
并取 $(\nu_1,\mu,\nu_0)\in\Theta(k)$。
将其在 $\{1,\ldots,k-1\}$ 上的限制
记为 $(\nu'_1,\mu',\nu'_0)\in\Theta(k-1)$。
\translatornote{原书式 (5.4.55)--(5.4.56)
把 $\mathfrak n$ 的右输入数与能量序号
写反；这里依
$\mathfrak n_{k_1,k_0;i}$ 的定义
和图 5.4.3 的中央顶点改为
$(\nu_1(k),\nu_0(k);\mu(k))$，
见疑误 A-0334。}
用 $x$、$z$ 的去拼接分解
$\Delta x=\sum_a x_a^{(1)}\otimes x_a^{(2)}$、
$\Delta z=\sum_b z_b^{(1)}\otimes z_b^{(2)}$
定义
\begin{equation}
\begin{aligned}
&\varphi_{(\nu_1,\mu,\nu_0)}(x\otimes y\otimes z)\\
&\quad=\sum_{a,b}\mathfrak G\circ
\mathfrak n_{\nu_1(k),\nu_0(k);\mu(k)}
\bigl(x_a^{(1)}\otimes
\varphi_{(\nu'_1,\mu',\nu'_0)}
(x_a^{(2)}\otimes y\otimes z_b^{(1)})
\otimes z_b^{(2)}\bigr),
\end{aligned}
\tag{5.4.55}\label{eq:5-4-55}
\end{equation}
以及
\begin{equation}
\begin{aligned}
&\mathfrak n'_{(\nu_1,\mu,\nu_0)}
(x\otimes y\otimes z)\\
&\quad=\sum_{a,b}\Pi_H\circ
\mathfrak n_{\nu_1(k),\nu_0(k);\mu(k)}
\bigl(x_a^{(1)}\otimes
\varphi_{(\nu'_1,\mu',\nu'_0)}
(x_a^{(2)}\otimes y\otimes z_b^{(1)})
\otimes z_b^{(2)}\bigr).
\end{aligned}
\tag{5.4.56}\label{eq:5-4-56}
\end{equation}
\translatornote{原书说 $\mathfrak G$、
$\Pi_H$ 分别延拓到 $\overline C$
与 $\overline D$；二者均是本证明中
$\overline D$ 上的收缩数据，
故都延拓到后者，见疑误 A-0328。}
这里 $\mathfrak G$ 与 $\Pi_H$
均平凡延拓到
$\overline D\otimes R[e,e^{-1}]$，
如定理 5.4.2 的证明中所做的那样。
\translatornote{原书下例把函数三元组
写成 $(3,\mu,1)$，而图示的左右输入数
实际为 $1$、$3$；
并把最下顶点的能量序号写成 $1$，
与图中 $\mu(1)=2$ 不符。
以下沿用已列明的三元组函数，
并在内层使用序号 $2$；
见疑误 A-0335、A-0336。}
例如，图 5.4.3 对应于
\[
\begin{aligned}
&\varphi_{(\nu_1,\mu,\nu_0)}
(x\otimes y\otimes z_1\otimes z_2\otimes z_3)\\
&\quad=\mathfrak G\circ\mathfrak n_{0,0;2}
\bigl(\mathfrak G\circ\mathfrak n_{1,2;0}
(x\otimes\mathfrak n_{0,1;2}(y\otimes z_1)
\otimes z_2\otimes z_3)\bigr),\\
&\mathfrak n'_{(\nu_1,\mu,\nu_0)}
(x\otimes y\otimes z_1\otimes z_2\otimes z_3)\\
&\quad=\Pi_H\circ\mathfrak n_{0,0;2}
\bigl(\mathfrak G\circ\mathfrak n_{1,2;0}
(x\otimes\mathfrak n_{0,1;2}(y\otimes z_1)
\otimes z_2\otimes z_3)\bigr).
\end{aligned}
\]
于是
\[
\begin{aligned}
\mathfrak n'_{k_1,k_0}
&=\sum_{k\ge0}
\sum_{\substack{(\nu_1,\mu,\nu_0)\in\Theta(k)\\
|\nu_1|=k_1,\ |\nu_0|=k_0}}
T^{|\mu|}\mathfrak n'_{(\nu_1,\mu,\nu_0)},\\
\varphi_{k_1,k_0}
&=\sum_{k\ge0}
\sum_{\substack{(\nu_1,\mu,\nu_0)\in\Theta(k)\\
|\nu_1|=k_1,\ |\nu_0|=k_0}}
T^{|\mu|}\varphi_{(\nu_1,\mu,\nu_0)}.
\end{aligned}
\]
按定理 5.4.2 证明中的同样方法，
可核验 $\mathfrak n'_{k_1,k_0}$
定义滤过 $A_\infty$ 双模结构，
而 $\varphi_{k_1,k_0}$
定义从 $D'$ 到 $D$ 的同伦等价。
此外，$D'$ 满足典范模型的全部性质。
若 $D$（同伦）含幺，
则可仿照定理 5.4.2$'$ 的证明
说明其典范模型也含幺。
定理 5.4.18 得证。
\end{proof}
